\documentclass{term-paper}

\title{Goodhart's Law and the Metrics We Choose}
\author{S. Lindqvist}
\course{Philosophy of Measurement}
\due{14 December}

\usepackage[style=authoryear,backend=bibtex]{biblatex}
\addbibresource{refs.bib}

\begin{document}

When a measure becomes a target, it ceases to be a good measure. This
is Goodhart's Law, and it is usually cited as a warning about
bureaucratic incentive design \parencite{strathern1997}. I argue here
that the law is not merely a caution about specific metrics gone
wrong, but a structural claim about any system where an agent both
optimises for a measure and controls the process the measure
observes --- and that this structural reading explains why the
failure recurs across domains that share almost nothing else in
common.

\section{The standard reading}

The standard reading of Goodhart's Law treats it as a story about
gaming: a call-centre worker rewarded for call volume rushes calls;
a school rewarded for test scores teaches to the test; a hospital
rewarded for short waiting times reclassifies patients rather than
treating them faster. In each case, the metric was a reasonable
\emph{proxy} for the underlying goal until it became the target
itself, at which point the agent's optimisation pressure found the
cheapest path to a high score rather than the intended path to the
underlying goal.

This reading is correct as far as it goes, but it treats each
instance as a separate failure of measurement design --- a
particular metric that turned out, in hindsight, to be gameable. That
framing invites a particular kind of response: design a better
metric, close the loophole, add a second metric to catch what the
first one missed. I want to suggest this response mistakes the
depth of the problem.

\section{A structural claim}

Consider what all three examples above share, beneath their surface
differences. In each, the agent being measured has some control over
the \emph{process} that produces the outcome the metric observes,
and that control is exercised with knowledge of the metric. This is
not a fact about call centres, schools or hospitals specifically ---
it is a fact about any principal-agent relationship where the agent
has more granular knowledge of the process than the principal has of
the outcome.

Under this reading, adding a second metric does not close the gap;
it relocates it. If a school is now measured on test scores
\emph{and} teacher-reported student wellbeing, the optimisation
pressure simply moves to whichever of the two measures has more slack
in it --- and since the agent, not the principal, has the finer-grained
knowledge of where that slack is, the principal cannot know in
advance which measure will absorb the pressure. This matches an
observation from complex-systems theory: a control system with fewer
degrees of freedom than the system it is trying to control cannot,
even in principle, capture that system's full behaviour
\parencite{ashby1956}. Two metrics do not have more degrees of freedom
than the one process generating both, if the agent understands the
relationship between them.

\section{Why this matters for design}

If the structural reading is right, the practical implication is not
``design better metrics'' but ``reduce the gap between what is
measured and what the agent actually controls, or reduce how much the
agent knows about what is measured.'' Both are harder and less
comfortable prescriptions than adding a second KPI, which may explain
why organisations reliably choose the second KPI instead: it produces
a defensible-looking response to a failure without requiring anyone
to give up the legibility that having a single clean number provides.

This also predicts something the gaming-story version does not: that
Goodhart failures should recur even in organisations that have
already been burned by one and have explicitly tried to correct for
it, because the correction (a second or third metric) does not close
the structural gap the argument above identifies. Anecdotally, this
matches what is actually observed --- organisations rarely escape
metric-gaming by adding metrics, and more often escape it (when they
do) by reducing an agent's discretion over the process itself, which
is a much more disruptive change than adjusting a scorecard.

\section{Conclusion}

Goodhart's Law is usually taught as a lesson about picking better
proxies. I have argued it is better understood as a claim about the
limits of any measurement relationship in which the measured party
has more knowledge of the underlying process than the measuring party
has of the outcome. This reframing does not make the problem easier
to solve, but it does explain why the easy solution --- adding another
metric --- keeps failing to solve it.

\printbibliography

\end{document}
